% !TEX program = xelatex
% Самодостаточный конспект. Компилировать XeLaTeX/LuaLaTeX (tectonic 03-irracionalnosti-spec-metody.tex).
\documentclass[11pt,a4paper]{article}
\newcommand{\coursename}{Математический анализ}
\newcommand{\rightheader}{§3. Иррациональности и специальные методы}
\newcommand{\doctitle}{Иррациональности и специальные методы интегрирования}
% ==== Общая преамбула конспектов (собирается в каждый .tex как самодостаточная) ====
% Перед этим файлом должны быть определены: \documentclass и макросы
%   \coursename  — название курса (левый колонтитул, подзаголовок),
%   \rightheader — правый колонтитул (напр. «§2. Рациональные дроби»),
%   \doctitle    — заголовок раздела на титуле.
\usepackage{fontspec}
\setmainfont{cmunrm.otf}[
  BoldFont       = cmunbx.otf,
  ItalicFont     = cmunti.otf,
  BoldItalicFont = cmunbi.otf]
\usepackage{polyglossia}
\setdefaultlanguage{russian}
\setotherlanguage{english}

\usepackage{amsmath,amssymb,amsthm,mathtools}
\usepackage[a4paper,margin=2.2cm]{geometry}
\usepackage{enumitem}
\usepackage{graphicx}
\usepackage{array}
\usepackage{xcolor}
\definecolor{accent}{HTML}{1F6FEB}
\definecolor{rulecol}{HTML}{D0D7DE}
\usepackage[colorlinks=true,linkcolor=accent,urlcolor=accent,citecolor=accent]{hyperref}
\usepackage{titlesec}
\usepackage{fancyhdr}

% ----- Колонтитулы -----
\pagestyle{fancy}
\fancyhf{}
\fancyhead[L]{\small\itshape \coursename}
\fancyhead[R]{\small\itshape \rightheader}
\fancyfoot[C]{\thepage}
\renewcommand{\headrulewidth}{0.4pt}
\renewcommand{\headrule}{\hbox to\headwidth{\color{rulecol}\leaders\hrule height \headrulewidth\hfill}}

% ----- Заголовки -----
\titleformat{\section}{\large\bfseries\color{accent}}{\thesection.}{0.6em}{}
\titleformat{\subsection}{\bfseries}{\thesubsection.}{0.5em}{}

% ----- Теоремные окружения (стиль конспекта) -----
\theoremstyle{definition}
\newtheorem{definition}{Определение}[section]
\newtheorem{example}[definition]{Пример}
\newtheorem{remark}[definition]{Замечание}
\theoremstyle{plain}
\newtheorem{theorem}[definition]{Теорема}
\newtheorem{lemma}[definition]{Лемма}
\newtheorem{corollary}[definition]{Следствие}
\newtheorem{property}[definition]{Свойство}
\newtheorem{criterion}[definition]{Критерий}
\renewcommand{\proofname}{Доказательство}

% ----- Операторы и сокращения -----
\DeclareMathOperator{\tg}{tg}
\DeclareMathOperator{\ctg}{ctg}
\DeclareMathOperator{\arctg}{arctg}
\DeclareMathOperator{\arcctg}{arcctg}
\DeclareMathOperator{\sh}{sh}
\DeclareMathOperator{\ch}{ch}
\DeclareMathOperator{\Th}{th}
\DeclareMathOperator{\cth}{cth}
\DeclareMathOperator{\Const}{const}
\DeclareMathOperator{\sign}{sign}
\DeclareMathOperator{\rang}{rang}
\DeclareMathOperator{\diam}{diam}
\DeclareMathOperator{\Span}{span}
\DeclareMathOperator{\Ker}{Ker}
\DeclareMathOperator{\Img}{Im}
\DeclareMathOperator{\tr}{tr}
\newcommand{\R}{\mathbb{R}}
\newcommand{\C}{\mathbb{C}}
\newcommand{\N}{\mathbb{N}}
\newcommand{\Z}{\mathbb{Z}}
\newcommand{\Q}{\mathbb{Q}}
\newcommand{\dx}{\,dx}
\newcommand{\dt}{\,dt}
\newcommand{\du}{\,du}
\newcommand{\eps}{\varepsilon}
\renewcommand{\Re}{\operatorname{Re}}
\renewcommand{\Im}{\operatorname{Im}}
\renewcommand{\le}{\leqslant}
\renewcommand{\ge}{\geqslant}
\newcommand{\scal}[2]{\left( #1,\,#2\right)}
\DeclarePairedDelimiter{\abs}{\lvert}{\rvert}
\DeclarePairedDelimiter{\norm}{\lVert}{\rVert}

\begin{document}
\begin{center}
  {\LARGE\bfseries \doctitle}\\[4pt]
  {\large\itshape \coursename}
\end{center}
\vspace{6pt}

\section{Дробно-линейные иррациональности}

Интеграл вида $\displaystyle\int R\!\left(x,\ \sqrt[n]{\dfrac{ax+b}{cx+d}}\right)\dx$ рационализуется
подстановкой
\[
  t=\sqrt[n]{\frac{ax+b}{cx+d}} .
\]
После неё $x$ выражается рационально через $t$, и подынтегральная функция становится рациональной.

\begin{example}
$\displaystyle\int\frac{1}{(1+x)^2}\sqrt{\frac{1-x}{1+x}}\dx$. Берём $t=\sqrt{\dfrac{1-x}{1+x}}$, тогда
$1-x=t^2(1+x)$, $x=\dfrac{1-t^2}{1+t^2}$, $1+x=\dfrac{2}{1+t^2}$, $\dx=-\dfrac{4t}{(1+t^2)^2}\dt$:
\[
  \int\frac{1}{(1+x)^2}\sqrt{\frac{1-x}{1+x}}\dx
   =\int\frac{(1+t^2)^2}{4}\cdot t\cdot\Bigl(-\frac{4t}{(1+t^2)^2}\Bigr)\dt
   =-\int t^2\dt=-\frac{t^3}{3}+C=-\frac13\Bigl(\frac{1-x}{1+x}\Bigr)^{3/2}+C.
\]
\end{example}

\section{Квадратичные иррациональности и подстановки Эйлера}

Для $\displaystyle\int R\bigl(x,\sqrt{ax^2+bx+c}\bigr)\dx$ применяют одну из трёх подстановок Эйлера,
каждая из которых рационализует интеграл.
\begin{enumerate}[label=\arabic*),leftmargin=*]
  \item \textbf{Первая} (при $a>0$): $\ \sqrt{ax^2+bx+c}=\pm\sqrt{a}\,x+t$.
  \item \textbf{Вторая} (при $c>0$): $\ \sqrt{ax^2+bx+c}=xt\pm\sqrt{c}$.
  \item \textbf{Третья} (при вещественных корнях $x_1,x_2$, т.е. $ax^2+bx+c=a(x-x_1)(x-x_2)$):
        $\ \sqrt{a(x-x_1)(x-x_2)}=t(x-x_1)$.
\end{enumerate}
В каждом случае возведение в квадрат даёт линейное относительно $x$ уравнение, откуда $x$ и
$\sqrt{ax^2+bx+c}$ выражаются рационально через $t$.

\begin{example}[третья подстановка]
$\displaystyle\int\frac{dx}{\sqrt{x^2-1}}$, где $x^2-1=(x-1)(x+1)$. Полагаем
$\sqrt{(x-1)(x+1)}=t(x-1)$, тогда $x+1=t^2(x-1)$, $x=\dfrac{t^2+1}{t^2-1}$,
$\dx=\dfrac{-4t}{(t^2-1)^2}\dt$, $x-1=\dfrac{2}{t^2-1}$, и интеграл сводится к рациональному; в
итоге $\displaystyle\int\frac{dx}{\sqrt{x^2-1}}=\ln\abs[\big]{x+\sqrt{x^2-1}}+C$ (см. также §3).
\end{example}

\section{Гиперболические и тригонометрические подстановки}

Для радикалов специального вида удобны подстановки, обращающие подкоренное выражение в полный
квадрат:
\[
  \sqrt{a^2-x^2}:\ x=a\sin t;\qquad
  \sqrt{a^2+x^2}:\ x=a\sh t;\qquad
  \sqrt{x^2-a^2}:\ x=a\ch t .
\]
Гиперболические подстановки часто предпочтительнее тригонометрических: они не требуют следить за
знаком и сразу приводят к интегралам от $\sh,\ch$, использующих $\ch^2 t-\sh^2 t=1$.

\begin{example}
$\displaystyle\int\sqrt{a^2+x^2}\dx$. Берём $x=a\sh t$, $\dx=a\ch t\dt$, $\sqrt{a^2+x^2}=a\ch t$:
\[
  \int\sqrt{a^2+x^2}\dx=a^2\!\int\ch^2 t\dt=a^2\!\int\frac{\ch 2t+1}{2}\dt
   =\frac{a^2}{2}\Bigl(\frac{\sh 2t}{2}+t\Bigr)+C,\qquad t=\operatorname{arsh}\frac xa .
\]
\end{example}

\begin{example}[вывод «длинного» логарифма]
$\displaystyle\int\frac{dx}{\sqrt{x^2-1}}$. Берём $x=\ch t$, $\dx=\sh t\dt$,
$\sqrt{x^2-1}=\sh t$:
\[
  \int\frac{dx}{\sqrt{x^2-1}}=\int\frac{\sh t\dt}{\sh t}=t+C=\operatorname{arch} x+C .
\]
Выразим $\operatorname{arch} x$ явно. Из $x=\ch t=\dfrac{e^t+e^{-t}}2$ для $y=e^t$ получаем
$y^2-2xy+1=0$, откуда $y=x+\sqrt{x^2-1}$ (берём знак $+$, так как $t\ge0$), и
\[
  \operatorname{arch} x=t=\ln\bigl(x+\sqrt{x^2-1}\bigr).
\]
\end{example}

\section{Интеграл от дифференциального бинома}

Дифференциальным биномом называют выражение $x^m(a+bx^n)^p$, где $a,b\in\R$ и $m,n,p\in\Q$.

\begin{theorem}[Чебышёв]
Интеграл $\displaystyle\int x^m(a+bx^n)^p\dx$ выражается через элементарные функции \emph{только} в
трёх случаях:
\begin{enumerate}[label=\arabic*),leftmargin=*]
  \item $p\in\Z$ — подстановка $x=t^k$, где $k$ — общий знаменатель дробей $m$ и $n$;
  \item $\dfrac{m+1}{n}\in\Z$ — подстановка $a+bx^n=t^k$, где $k$ — знаменатель $p$;
  \item $\dfrac{m+1}{n}+p\in\Z$ — подстановка $ax^{-n}+b=t^k$, где $k$ — знаменатель $p$.
\end{enumerate}
Иначе интеграл не выражается в элементарных функциях.
\end{theorem}

\begin{example}
$\displaystyle\int\frac{\sqrt[3]{1+\sqrt[4]{x}}}{\sqrt{x}}\dx=\int x^{-1/2}\bigl(1+x^{1/4}\bigr)^{1/3}\dx$.
Здесь $m=-\tfrac12$, $n=\tfrac14$, $p=\tfrac13$, и $\dfrac{m+1}{n}=2\in\Z$ — второй случай. Замена
$1+x^{1/4}=t^3$ даёт
\[
  \int\frac{\sqrt[3]{1+\sqrt[4]{x}}}{\sqrt{x}}\dx
   =12\int(t^6-t^3)\dt=\frac{12}{7}\bigl(1+\sqrt[4]{x}\bigr)^{7/3}-3\bigl(1+\sqrt[4]{x}\bigr)^{4/3}+C.
\]
\end{example}

\begin{remark}
Например, $\displaystyle\int\sqrt{1+x^3}\dx$ \emph{не берётся}: здесь $m=0$, $n=3$, $p=\tfrac12$, и
ни одно из трёх условий Чебышёва не выполнено.
\end{remark}

\section{Параметрическое дифференцирование}

Правило Лейбница (дифференцирование интеграла по параметру) позволяет порождать новые интегральные
формулы из уже известных, дифференцируя обе части по параметру.

\begin{example}
Из $\displaystyle\int\cos(\alpha x)\dx=\frac{\sin(\alpha x)}{\alpha}+C$ дифференцированием по
$\alpha$ получаем $\displaystyle\int -x\sin(\alpha x)\dx=\frac{x\cos(\alpha x)}{\alpha}
-\frac{\sin(\alpha x)}{\alpha^2}$, то есть
\[
  \int x\sin(\alpha x)\dx=\frac{\sin(\alpha x)}{\alpha^2}-\frac{x\cos(\alpha x)}{\alpha}+C.
\]
\end{example}

\section{Классификация «неберущихся» интегралов}

Не всякая элементарная функция имеет элементарную первообразную. Классические примеры
интегралов, \emph{не выражающихся} через элементарные функции:
\begin{itemize}[leftmargin=*]
  \item \textbf{Интеграл Гаусса (Пуассона):} $\displaystyle\int e^{-x^2}\dx$ — задаёт функцию
        ошибок $\operatorname{erf}(x)=\dfrac{2}{\sqrt\pi}\displaystyle\int_0^x e^{-t^2}\dt$.
  \item \textbf{Интегральный синус и косинус:} $\displaystyle\operatorname{Si}(x)=\int_0^x\frac{\sin t}{t}\dt$,
        $\displaystyle\operatorname{Ci}(x)$.
  \item \textbf{Эллиптические интегралы:} $\displaystyle\int\frac{dx}{\sqrt{1-k^2\sin^2 x}}$ и
        родственные — возникают, в частности, при вычислении длины дуги эллипса.
  \item \textbf{Случаи Чебышёва:} дифференциальные биномы, не подпадающие под три случая теоремы
        (например $\int\sqrt{1+x^3}\dx$).
\end{itemize}
Такие интегралы определяют новые \emph{специальные функции}, изучаемые отдельно.

\end{document}
